\documentclass[11pt]{article}
\usepackage[margin=1in]{geometry}
\usepackage{amsmath,amssymb,amsthm,microtype}
\usepackage[colorlinks=true,urlcolor=blue,citecolor=blue]{hyperref}
\newcommand{\Mag}{\operatorname{Mag}}
\newcommand{\ellone}{\ell_1}

\title{Compact positive-definite spaces can have infinite magnitude}
\author{Literature-status note for \texttt{fa\_banach\_001}}
\date{27 August 2026}

\begin{document}
\maketitle

\noindent\textbf{Status.}
\texttt{literature\_already\_answered (negative answer)}.  No new
mathematical result is claimed.

\section*{Original question}

For a compact positive-definite metric space $X$, Meckes defines magnitude by
the supremum of the finite quadratic ratios in display (1) of
\cite{Meckes2019}.  Immediately after that display, on arXiv PDF page 2, he
asks whether this supremum is finite for every compact positive-definite
metric space.

\section*{Explicit later answer}

Leinster and Meckes explicitly identify this question on page 2 of
\cite{LeinsterMeckes2021}.  They state that no compact positive-definite
space of infinite magnitude had previously been known, cite
\cite{Meckes2019} among the papers in which the question was raised, and
announce examples with infinite magnitude at every positive scale.

Their Theorem 2.1, on arXiv PDF page 3, gives the answer.  If $(a_i)$ is a
positive null sequence with $\sum_i a_i=\infty$, and $X$ is the closed convex
hull of $\{a_i e_i:i\geq 1\}$ in $\ellone$, then $X$ is compact and of
negative type, and
\[
  \Mag(tX)=\infty \qquad (t>0).
\]
Negative type means precisely that every positive rescaling is positive
definite.  In particular, $X$ itself is compact and positive definite, so the
theorem is an exact negative answer to the source question, with the stronger
all-scales conclusion.

\section*{A short star-subspace verification}

There is also a direct nonconvex witness inside the same ambient space.  Put
\[
  S=\{0\}\cup\{x_n:n\geq 1\}\subset\ellone,
  \qquad x_n=\frac1n e_n.
\]
Then $x_n\to0$, so $S$ is compact.  Since $\ellone$ is of negative type,
$S$ is positive definite.  For
$S_N=\{0,x_1,\ldots,x_N\}$, write $a_i=1/i$ and $q_i=e^{-a_i}$.  The
similarity matrix $Z=(e^{-d(x,y)})_{x,y\in S_N}$ satisfies
\[
 Z_{0i}=q_i,\qquad Z_{ij}=q_iq_j\quad(i\ne j),\qquad Z_{ii}=1.
\]
Solving $Zw=\mathbf 1$, set
$T=w_0+\sum_{j=1}^Nq_jw_j$.  The zeroth row gives $T=1$, while row $i$
gives
\[
 q_iT+(1-q_i^2)w_i=1,
 \qquad\text{hence}\qquad w_i=\frac1{1+q_i}.
\]
Consequently,
\[
 \Mag(S_N)=\sum_{i=0}^N w_i
 =1+\sum_{i=1}^N\frac{1-q_i}{1+q_i}
 =1+\sum_{i=1}^N\tanh\!\left(\frac1{2i}\right).
\]
The last sum diverges with $N$.  Since compact positive-definite magnitude is
the supremum of the magnitudes of finite subspaces, $\Mag(S)=\infty$.
Replacing $1/i$ by $t/i$ proves $\Mag(tS)=\infty$ for every $t>0$.

This elementary calculation independently verifies the negative answer, but
it does not change the provenance classification: the exact question was
already answered explicitly by \cite{LeinsterMeckes2021}.

\section*{Scope and search status}

The identification was checked in the cached arXiv sources and PDFs.  The
source question is on page 2 of arXiv:1904.08923.  The answering paper names
the prior question on page 2 and proves its stronger Theorem 2.1 on page 3.
The run's cheap indexes also contain a prior triage record for arXiv:2112.12889;
that record correctly treats the later paper's own introductory signal as a
same-paper ask-and-answer extraction.  It does not negate the two-paper
literature relation recorded here.

Other conjectures in arXiv:1904.08923, including the small-scale derivative
formula for convex bodies and the Lipschitz-distance continuity conjecture,
are outside this packet and are not claimed resolved.

\begin{thebibliography}{9}
\bibitem{Meckes2019}
M.~W. Meckes,
\emph{On the magnitude and intrinsic volumes of a convex body in Euclidean
space}, Mathematika \textbf{66} (2020), 343--355;
arXiv:1904.08923,
\url{https://arxiv.org/abs/1904.08923}.

\bibitem{LeinsterMeckes2021}
T.~Leinster and M.~W. Meckes,
\emph{Spaces of extremal magnitude},
arXiv:2112.12889,
\url{https://arxiv.org/abs/2112.12889}.
\end{thebibliography}

\end{document}
